\section*{Chapter \ref{ch:rc:convexity-adjustments-and-constant-maturity-products} --- Convexity Adjustments and Constant-Maturity Products}

\begin{solution}{exo:rc:convexity-adjustments-and-constant-maturity-products:1}
$\delta\sigma_N^2T/(1+\delta F) = 0.25\times0.008^2\times4/(1+0.25\times0.025) =
0.64$ basis points.
\end{solution}

\begin{solution}{exo:rc:convexity-adjustments-and-constant-maturity-products:2}
Paid soon after observation, the payment measure's numeraire $P(T,T_p)$ is
nearly one while the annuity falls when rates rise, so $P/A$ rises with $S$:
$a_1>0$, and the payment measure weights high-rate states more than the
annuity measure does. With positive rates $a_0 + a_1S(0) = P(0,T_p)/A(0) >
1/\sum\delta_i = a_0$, hence $a_1 > 0$.
\end{solution}

\begin{solution}{exo:rc:convexity-adjustments-and-constant-maturity-products:3}
$a_0 = 1/(10\times1.0139) = 0.09863$.
\end{solution}

\begin{solution}{exo:rc:convexity-adjustments-and-constant-maturity-products:4}
With a flat smile the adjustment is proportional to $\sigma_N^2$:
$11.73\times1.2^2 = 16.89$ basis points.
\end{solution}

\begin{solution}{exo:rc:convexity-adjustments-and-constant-maturity-products:5}
About 29\%: swaptions more than two percentage points out of the money carry
nearly a third of the ten-year adjustment. The cube's wings, which the swaption
market quotes thinly, must be marked with care, bounded by densities
(chapter~5), and are often calibrated to traded CMS spreads.
\end{solution}

\begin{solution}{exo:rc:convexity-adjustments-and-constant-maturity-products:6}
$+9$ basis points: when euro rates rise the euro tends to fall, so each euro basis
point is worth fewer dollars in the states where rates are high; the dollar
payment measure shifts the rate's expectation up.
\end{solution}

\begin{solution}{exo:rc:convexity-adjustments-and-constant-maturity-products:7}
$-19.24$ basis points. Paid at the swap's end, $P(T,T_p)/A(T)$ falls when rates
rise (the long discount factor falls faster than the annuity), so $a_1<0$ and
the payment measure weights low-rate states more: the adjustment is negative.
\end{solution}

\begin{solution}{exo:rc:convexity-adjustments-and-constant-maturity-products:8}
Each CMS coupon is a swap rate paid once, on a date that does not match the
annuity numeraire under which the forward swap rate is a martingale; its value
is the forward plus a convexity adjustment worth a strip of swaptions (32
basis points for the ten-year rate in ten years here). Valued at forwards the
coupons are underpriced, and a hedge with forward swaps leaves the book short
volatility; hedge the adjustment with swaptions across strikes.
\end{solution}

\begin{solution}{pb:rc:convexity-adjustments-and-constant-maturity-products:1}
\textbf{1.} A \omterm{def:rc:convexity-adjustments-and-constant-maturity-products:spread}{CMS spread option} struck at zero on the ten-year minus two-year rate,
paid a year after observation, times $p$.
\textbf{2.} Each is a swap rate paid once, a year after observation: each needs a
CMS adjustment; the ten-year's is larger, its annuity being more sensitive to
rates and its volatility higher per unit of $a_1$.
\textbf{3.} 0.17118 per unit notional.
\textbf{4.} 0.05785.
\textbf{5.} $0.17118/0.05785 = 2.96$.
\textbf{6.} 2.44, 3.10 and 4.47.
\textbf{7.} Higher correlation, lower spread volatility, cheaper spread floors
struck at zero (the spread's forward being positive): each unit of coupon
costs less, so more units fit in the same budget.
\textbf{8.} It is short the spread options it pays, so short spread volatility and
long correlation: it loses if correlation falls.
\textbf{9.} 45.1 basis points.
\textbf{10.} Swap rates and Treasury yields differ by swap spreads, which move with
their own dynamics; the correlation that matters is under the pricing measure
and over the note's horizon, not daily historical changes.
\textbf{11.} Out-of-the-money payer and receiver swaptions on each rate's
underlying swap, across strikes, in the replication weights.
\textbf{12.} The CMS adjustments rise, raising the CMS forwards of both rates; the
effect on the spread depends on which rises more, and the spread options gain
from the higher volatility of each leg.
\textbf{13.} Both: the level through the CMS adjustments and the annuity mapping,
the slope directly through the spread's forward.
\textbf{14.} Coupons fall to zero: the floor at zero protects the investor from
negative coupons, not from receiving nothing for years.
\textbf{15.} Many investors buy the same trade from few dealers, who end up short
the same correlation and spread volatility and can hedge it only with each
other; the book grows with issuance and is hard to reduce.
\textbf{16.} Lower than 2.96: the bank charges for its hedging costs, correlation
reserve and funding; a margin of a few tenths of participation is usual for
illiquid correlation.
\textbf{17.} From historical estimates of several windows and regimes, from the
correlations implied by any traded spread options, with a reserve for the
spread between them; stress it in the risk report.
\textbf{18.} That the note pays only if the curve is steeper than zero at each
fixing; a flat or inverted curve for years pays nothing, and the forwards
already price the expected steepening.
\textbf{19.} Named result: \emph{the steepener note} can pay 2.96 times the ten-year
minus two-year spread at a correlation of 0.769, between 2.44 and 4.47 as the
correlation goes from 0.6 to 0.95.
\textbf{20.} A strip of options on the curve's slope, paid for by giving up a fixed
coupon.
\end{solution}

\begin{solution}{iq:rc:convexity-adjustments-and-constant-maturity-products:1}
The difference between a rate's expectation under the measure of its actual
payment and its forward, caused by the correlation between the rate and the
ratio of numeraires. Examples: futures against forwards (daily margining),
a rate paid at its fixing date (timing), a CMS rate, a rate paid in another
currency (quanto).
\iqlookfor{the measure-change explanation and several instances.}
\end{solution}

\begin{solution}{iq:rc:convexity-adjustments-and-constant-maturity-products:2}
Change from the $T_k$- to the $T_{k-1}$-forward measure with density
$(1+\delta F_k(T_{k-1}))/(1+\delta F_k(0))$; then $\E^{T_{k-1}}[F] =
F_0+\delta\Var^{T_k}(F)/(1+\delta F_0)$, i.e. $\delta\sigma_N^2T/(1+\delta F_0)$ in the
normal model.
\iqlookfor{the Radon--Nikodym density between forward measures.}
\end{solution}

\begin{solution}{iq:rc:convexity-adjustments-and-constant-maturity-products:3}
Write the payoff under the annuity measure: value $=A(0)\E^A[(S-K)^+P(T,T_p)/A(T)]$;
map $P/A$ to $a_0+a_1S$ (\omterm{def:rc:convexity-adjustments-and-constant-maturity-products:tsr}{linear terminal swap-rate model}); expand
$(S-K)^+(a_0+a_1S) = (a_0+a_1K)(S-K)^+ + 2a_1\int_K^\infty(S-L)^+dL$ and price each
payer swaption on the smile.
\iqlookfor{annuity mapping plus static replication.}
\end{solution}

\begin{solution}{iq:rc:convexity-adjustments-and-constant-maturity-products:4}
Down: the CMS coupons you owe are worth the forward plus a strip of
out-of-the-money swaptions, so you are short volatility, including the wings;
strikes around and well away from the forward matter, nearly a third of the
adjustment coming from more than two percentage points away.
\iqlookfor{short vega through the replication, and the wings.}
\end{solution}

\begin{solution}{iq:rc:convexity-adjustments-and-constant-maturity-products:5}
Mark from a range of historical estimates and any implied correlations from
traded spread options or dealer polls, with a reserve for the uncertainty;
hedge the correlation exposure with offsetting spread products where
possible, otherwise run it within limits and stress it; report the correlation
sensitivity separately.
\iqlookfor{marks with reserves, and honesty about unhedgeable risk.}
\end{solution}

\begin{solution}{iq:rc:convexity-adjustments-and-constant-maturity-products:6}
Grid width (enough standard deviations, and no strikes below the model's floor),
spacing near the forward where the integrand is steep, whether the strike and
the forward are grid nodes, the smile's behaviour in the wings (extrapolation,
NaNs), and convergence as the grid is refined; compare with the closed form
under a flat smile.
\iqlookfor{numerical integration hygiene and a benchmark.}
\end{solution}
